\section{Sparse and Efficient Attention Mechanisms}
\label{sec:sparse}

The quadratic complexity of standard self-attention mechanisms represents a fundamental bottleneck for processing long sequences in transformer-based language models. Standard self-attention scales as $O(n^2)$ with respect to sequence length $n$, imposing severe constraints on both computational resources and memory capacity. This quadratic scaling creates a practical ceiling on context length, typically limiting transformers to sequences of 512 to 2,048 tokens. As the demand for longer context windows has grown across applications ranging from document analysis to code understanding, researchers have pursued diverse strategies to mitigate this fundamental limitation. The approaches fall into several broad categories: fixed sparse attention patterns that reduce the number of attended positions, linear approximations that reformulate the attention mechanism itself, state space models that replace attention with alternative sequence mixing operations, and hardware-aware optimizations that maintain quadratic complexity while dramatically improving wall-clock performance through algorithmic efficiency.

\subsection{Taxonomy of Efficient Attention}

\citet{tay2020efficient} provided a comprehensive survey of efficient transformer architectures, establishing a taxonomy that has guided subsequent research. Their analysis identified that attention mechanisms can be modified along multiple dimensions, including the sparsity pattern imposed on attention weights, the computational approach used to approximate or compute attention, and the architectural paradigm employed for sequence modeling. Sparsity-based methods reduce the number of token pairs that participate in attention computation, either through fixed patterns determined by position or through learned patterns based on content. Approximation methods maintain dense attention but reduce computational complexity through mathematical reformulations, including kernel methods and low-rank factorizations. Alternative architectures abandon the attention mechanism entirely, replacing it with state space models, linear recurrences, or hybrid approaches that combine multiple paradigms.

\subsection{Fixed Sparse Attention Patterns}

\subsubsection{Longformer}

\citet{beltagy2020longformer} introduced the Longformer architecture, which represents one of the most influential approaches to sparse attention through the combination of local windowed attention with task-motivated global attention. The local component employs sliding windows of fixed size $w$, typically set to 512 tokens, where each token attends to $w/2$ neighbors on either side. This windowed attention captures local dependencies and maintains fine-grained contextual information within the immediate vicinity of each token. The windowed approach can be extended with dilated windows, analogous to dilated convolutions in computer vision, enabling the model to capture longer-range patterns without incurring full quadratic cost.

The global attention component allows designated tokens to attend to all positions in the sequence and to receive attention from all other tokens. For classification tasks, the \texttt{[CLS]} token typically receives global attention, enabling it to aggregate information from the entire document. For question answering, both question tokens and special separator tokens may employ global attention to facilitate cross-document reasoning. This combination of local and global patterns achieves $O(n)$ complexity with respect to sequence length while maintaining expressiveness comparable to full attention.

Longformer demonstrated strong empirical results across multiple domains. On character-level language modeling benchmarks including text8 and enwik8, Longformer achieved state-of-the-art results, validating its ability to model long-range dependencies in sequential data. The architecture consistently outperformed RoBERTa on long document tasks, with particularly strong results on WikiHop question answering, achieving 65.10\% F1 score, and TriviaQA, reaching 80.51% F1. The model handles sequences from 1,024 to 4,096 tokens effectively, with computational advantages that scale linearly with sequence length. At a sequence length of 2,048 tokens, Longformer achieves approximately 10x memory savings compared to standard attention, with even greater benefits at 4,096 tokens where savings reach 20x. The architecture serves as a drop-in replacement for standard self-attention and integrates seamlessly into existing transformer models, with pretrained Longformer variants based on the RoBERTa architecture available through open-source implementations.

Beyond the core architecture, \citet{beltagy2020longformer} introduced the Longformer-Encoder-Decoder (LED) variant for sequence-to-sequence tasks, extending the sparse attention paradigm to encoder-decoder architectures. This extension enables efficient processing of long-document summarization and translation tasks where both input and output sequences may be lengthy. The practical impact of Longformer has been substantial, with widespread adoption in applications requiring long-context understanding including legal document analysis, scientific paper processing, and genomic sequence modeling.

\subsubsection{BigBird}

\citet{zaheer2020bigbird} proposed BigBird, a sparse attention mechanism that combines three complementary components to achieve linear complexity while preserving theoretical properties of full attention. The first component consists of global tokens, a fixed set of $g$ positions that attend to all other positions and receive attention from the entire sequence. These global tokens serve as information hubs, aggregating and distributing information across distant portions of the sequence. The second component employs local window attention, where each token attends to $w$ neighboring positions, capturing fine-grained local context similar to Longformer's windowed approach. The third and most distinctive component introduces random attention connections, where each query attends to $r$ randomly selected keys throughout the sequence.

The random attention component represents a key theoretical and practical contribution. Inspired by graph sparsification techniques, random connections provide probabilistic coverage of long-range dependencies without requiring the model to specify which distant positions are relevant a priori. This randomness introduces diversity in the attention pattern, enabling the discovery of unexpected dependencies that might be missed by purely local or global patterns. The combination of these three components yields an attention mechanism with $O(n)$ complexity while maintaining expressiveness.

BigBird's theoretical contributions extend beyond empirical performance. \citet{zaheer2020bigbird} proved that BigBird sparse attention is a universal approximator of sequence functions, demonstrating that the sparse pattern can approximate any function that full attention could compute, given sufficient model capacity. Furthermore, they established that BigBird is Turing complete, showing that the architecture preserves the computational expressiveness of standard transformers despite the sparsity constraint. These theoretical guarantees distinguish BigBird from heuristic sparsity patterns and provide confidence that the sparse structure does not fundamentally limit model capabilities.

Empirically, BigBird handled sequences up to 8x longer than standard transformers could process on equivalent hardware, effectively managing sequences of 4,096 tokens and demonstrating viability for sequences extending to 8,000 tokens. The architecture achieved significant improvements on long-document question answering tasks including NaturalQuestions and HotpotQA, where the ability to reason across lengthy contexts proved crucial. Beyond natural language, BigBird demonstrated remarkable effectiveness on genomic sequence tasks, setting state-of-the-art results on DNA sequence modeling where sequences can extend to tens of thousands of base pairs. The successful application to genomics validated that sparse attention mechanisms generalize beyond language to other sequential domains with long-range dependencies.

The BigBird architecture encompasses multiple variants optimized for different use cases. BigBird-ITC employs internal transformer construction with sparse attention applied uniformly across layers, while BigBird-ETC uses extended transformer construction that varies the sparsity pattern across different layers of the network. These variants provide flexibility in trading off between computational efficiency and model expressiveness depending on task requirements.

\subsection{Linear Attention Approximations}

\subsubsection{Performer}

\citet{choromanski2021rethinking} introduced the Performer, which achieves linear-time attention through the FAVOR+ (Fast Attention Via positive Orthogonal Random features) algorithm. Unlike sparse attention methods that reduce the number of attended positions, Performer maintains dense attention but approximates the attention computation through kernel methods. The core insight is that softmax attention can be expressed as a kernel function, and kernel functions can be approximated using random feature maps. By decomposing the attention matrix using positive orthogonal random features, Performer avoids the explicit computation of the $n \times n$ attention matrix.

The FAVOR+ algorithm provides several theoretical guarantees that distinguish it from earlier linear attention approximations. The approach yields unbiased or nearly-unbiased estimates of the true attention weights, ensuring that the expected value of the approximation matches the exact computation. The estimation exhibits low variance, providing consistent results across different random seeds and training runs. Uniform convergence properties ensure that the approximation quality improves predictably as the number of random features increases. These mathematical guarantees provide confidence in the approximation quality and enable principled tuning of the feature dimension to balance accuracy and efficiency.

The Performer achieves $O(n)$ time complexity and $O(n)$ space complexity with respect to sequence length, compared to $O(n^2)$ for both in standard attention. The approach maintains compatibility with standard transformer architectures, requiring minimal modification to existing codebases. Beyond softmax attention, FAVOR+ supports kernelizable attention mechanisms more broadly, enabling exploration of alternative attention formulations including polynomial kernels and custom kernel functions designed for specific domains.

Empirically, Performer demonstrated competitive performance across multiple modalities. On language modeling tasks, the architecture achieved results comparable to standard transformers while enabling processing of longer sequences within fixed memory budgets. For pixel-level image prediction in the style of ImageGPT, Performer successfully modeled high-resolution images as long sequences of pixels. The architecture proved particularly effective on protein sequence modeling, where the ability to process full-length proteins as single sequences enabled improved structure prediction. On standard benchmarks including the Long Range Arena suite, Performer achieved competitive results while maintaining the theoretical linear complexity guarantee.

The practical speedup achieved by Performer depends on sequence length, with benefits increasing as sequences grow longer. For moderately long sequences of approximately 1,000 tokens, Performer achieves 3-4x speedup compared to standard attention implementations. The speedup increases substantially for sequences of several thousand tokens, where memory constraints often prevent standard attention from running at all. The random feature approximation introduces a small quality degradation compared to exact attention, typically on the order of 1-2\% on downstream metrics, representing a reasonable trade-off for applications where sequence length or computational efficiency are critical.

\subsubsection{Linformer}

\citet{wang2020linformer} demonstrated that the self-attention mechanism can be approximated by low-rank matrices, enabling a different path to linear complexity. The key theoretical insight is that attention matrices in trained transformers often exhibit approximately low-rank structure, suggesting that much of the information can be captured in a lower-dimensional subspace. Linformer exploits this structure by projecting the keys and values to a lower-dimensional space before computing attention, reducing the effective sequence length from $n$ to $k$ where $k \ll n$.

The projection strategy employs learned linear transformations that map the key and value matrices from dimension $n \times d$ to $k \times d$, where $k$ is the projection dimension and $d$ is the model dimension. By performing this projection, the subsequent attention computation operates on effective sequence length $k$ rather than $n$, yielding complexity $O(nk)$ which is linear in $n$ when $k$ is treated as a constant. The projection matrices can be shared across different layers or attention heads to reduce parameters, or they can be layer-specific and head-specific to maximize expressiveness. Position-aware projections that incorporate positional information into the projection operation can further improve performance.

With typical projection dimensions of $k=128$ or $k=256$ for sequences of length $n=512$, Linformer achieves substantial efficiency gains. The approach yields 1.5x faster inference time and enables 1.7x larger maximum batch sizes due to reduced memory consumption. These benefits scale favorably as sequences grow longer, with 10x memory reduction at sequence length 2,048 and 20x reduction at 4,096 tokens. Critically, Linformer maintains performance comparable to standard transformers on language modeling benchmarks, with minimal degradation in perplexity or downstream task accuracy.

The low-rank approximation approach provides interpretability benefits, as the learned projection matrices reveal which aspects of the sequence the model considers most important for attention computation. Analysis of trained Linformer models shows that projections often learn to emphasize content-bearing tokens while down-weighting function words and positional markers, suggesting that the approximation captures semantically meaningful structure. The approach proved effective across diverse applications including question answering on SQuAD, long document classification, and machine translation, demonstrating broad applicability of the low-rank attention paradigm.

\subsection{State Space Models}

\subsubsection{Structured State Spaces (S4)}

\citet{gu2022efficiently} introduced S4 (Structured State Space), representing a fundamental departure from attention-based architectures by modeling sequences through continuous-time state space equations. State space models describe dynamical systems through the equations $\frac{dx}{dt} = Ax + Bu$ and $y = Cx + Du$, where $x$ is a latent state, $u$ is the input, and $y$ is the output. These equations, foundational in control theory and signal processing, capture how systems evolve over time while maintaining a compressed representation of history in the state vector.

The key innovation in S4 is the structured parameterization of the state transition matrix $A$ through HiPPO (High-order Polynomial Projection Operators), which provide a principled approach to initializing the state dynamics. HiPPO matrices are designed to optimally compress history into a finite-dimensional state representation by projecting the history onto a basis of orthogonal polynomials. This initialization enables the model to capture long-range dependencies from the start of training rather than requiring the model to discover appropriate time scales through gradient descent.

S4 achieves efficient computation through a dual representation that expresses the same model as either a convolution or a recurrence. During training, the convolutional representation enables parallel processing of entire sequences using FFT-based algorithms, achieving $O(n \log n)$ complexity. During inference, the recurrent representation processes tokens sequentially with constant memory, requiring only the maintenance of a fixed-size state vector regardless of sequence length. This duality provides the best of both paradigms: parallelizable training akin to transformers and efficient constant-memory inference akin to RNNs.

The structured parameterization of $A$ proves critical for computational efficiency. By imposing structure through low-rank corrections to diagonal matrices, S4 enables stable diagonalization of the state matrix. This diagonalization reduces the state space computation to a Cauchy kernel, which admits closed-form evaluation with favorable numerical properties. The reduction to Cauchy kernels enables efficient and numerically stable computation even for sequences of tens of thousands of tokens, addressing a key weakness of earlier state space model approaches that suffered from numerical instability.

S4 achieved remarkable empirical results, particularly on tasks requiring long-range reasoning. The model became the first to solve the Long Range Arena Path-X task, which requires tracking dependencies across 16,384 positions, achieving 88\% accuracy compared to 50\% random baseline and near-zero performance from standard transformers. On Sequential CIFAR-10, where images are processed as sequences of 1,024 pixels, S4 achieved 91\% accuracy without data augmentation, demonstrating effectiveness on long visual sequences. For speech classification tasks with audio sequences of length 16,000, S4 achieved 1.7\% error rate, representing a 2x improvement over previous methods.

The architecture demonstrated 60x faster generation compared to standard transformers of equivalent size, enabled by the constant-memory recurrent inference mode. This efficiency advantage proved particularly valuable for applications requiring generation of long sequences, including audio synthesis and long-form text generation. S4's success on raw audio generation tasks validated its ability to capture structure across multiple time scales, a critical requirement for modeling acoustic signals with both fine-grained timing and long-term musical or linguistic structure.

\subsubsection{Mamba}

\citet{gu2023mamba} introduced Mamba, advancing state space models through selective state space mechanisms that address a key limitation of prior SSMs. While S4 and related models achieved impressive efficiency, they struggled on discrete, information-dense data like natural language where the ability to selectively filter and route information proves crucial. Mamba introduces selectivity by making SSM parameters—specifically the discretization parameter $\Delta$, input matrix $B$, and output matrix $C$—functions of the input content rather than fixed parameters.

This input-dependent parameterization enables content-based reasoning analogous to the query-key-value mechanism in attention. By allowing the model to dynamically adjust how information flows through the state based on token content, Mamba can selectively propagate relevant information and filter irrelevant content. This selectivity addresses the fundamental challenge that fixed-parameter SSMs face when processing language, where importance and relevance of tokens varies dramatically based on semantic content rather than following fixed patterns.

The selective mechanism is implemented through a hardware-efficient scan algorithm that computes state updates efficiently on GPUs. Following principles from FlashAttention, the implementation minimizes memory movement between GPU memory hierarchies, enabling the selective computation to run efficiently despite the additional complexity of input-dependent parameters. The architecture employs a simplified design compared to transformers, eliminating separate attention and MLP blocks in favor of a single Mamba block that combines state space dynamics with gating mechanisms.

Mamba achieves linear $O(n)$ time complexity while maintaining constant-space inference with no KV cache requirement. The throughput advantages are substantial: Mamba achieves 5x higher throughput than transformers of equivalent size, with throughput scaling linearly rather than quadratically as sequence length increases. The architecture handles sequences up to one million tokens with constant memory footprint, enabling applications previously infeasible with transformer-based models.

The empirical results validated Mamba's approach across multiple modalities. For language modeling, Mamba-3B outperformed transformers of the same size and matched the performance of transformers twice its size, demonstrating that the selective state space mechanism provides sufficient expressiveness for language tasks despite abandoning attention entirely. On downstream language tasks including question answering, classification, and reasoning benchmarks, Mamba exhibited strong zero-shot and few-shot capabilities comparable to attention-based models. Beyond language, Mamba achieved state-of-the-art results on audio modeling tasks, competitive performance on genomic sequence analysis, and promising results on image understanding when adapted to vision tasks.

The impact of Mamba has been substantial, accumulating over 5,000 citations by 2024 and sparking extensive follow-up research. The architecture demonstrated that alternatives to attention can achieve competitive quality on the full range of tasks that have driven transformer development, while offering substantial efficiency advantages. The success of Mamba has renewed interest in recurrent architectures and state space models as viable alternatives to attention for sequence modeling at scale.

\subsubsection{Hybrid Architectures: Jamba}

\citet{lieber2024jamba} introduced Jamba, the first production-grade hybrid architecture that combines Transformer attention layers, Mamba state space layers, and Mixture-of-Experts (MoE) into a unified decoder architecture. The hybrid design emerged from the observation that attention and state space models offer complementary strengths: attention excels at capturing fine-grained dependencies and content-based routing, while state space models provide efficient long-range modeling and constant-memory inference. By interleaving these components, Jamba aims to leverage the advantages of both paradigms while mitigating their respective weaknesses.

The Jamba architecture employs a typical ratio of one Transformer layer to seven Mamba layers, creating repeating blocks that combine both architectural paradigms. The Transformer layers provide global attention capacity at strategic intervals, enabling the model to perform explicit content-based retrieval and reasoning. The Mamba layers handle the bulk of sequence processing with linear complexity, maintaining long-range state while avoiding the quadratic cost of full attention. Selected layers incorporate Mixture-of-Experts, where input is routed to a subset of expert networks, enabling parameter efficiency through conditional computation. With 52B total parameters but only 12B active parameters per token, the MoE component provides substantial model capacity while maintaining computational efficiency.

The hybrid approach yields impressive practical benefits. Jamba achieves 3x higher throughput than pure transformers of equivalent size, enabled by the predominance of efficient Mamba layers. The architecture supports a 256K token context window, among the largest available at the time of release, demonstrating that hybrid designs can scale to extreme context lengths. The combination of linear-complexity Mamba layers and selective attention allows the model to process long documents efficiently while maintaining strong performance on tasks requiring precise attention to detail.

Empirically, Jamba demonstrated quality competitive with larger pure-transformer models, validating that the hybrid architecture does not sacrifice capability for efficiency. The model proved particularly effective on long-context tasks including document question answering, code understanding, and retrieval-augmented generation scenarios. The production-ready implementation and Apache 2.0 open license facilitated adoption, with Jamba serving as a proof of concept that hybrid architectures can succeed in practical deployment.

The Jamba architecture suggests that future language models may increasingly employ heterogeneous layer types, selecting architectural components based on their strengths for specific computational roles. Rather than committing entirely to attention or alternatives, hybrid approaches enable fine-grained optimization of the efficiency-expressiveness trade-off across different layers of deep networks.

\subsection{Linear-Time Alternatives}

\subsubsection{RWKV}

\citet{peng2023rwkv} proposed RWKV (Receptance Weighted Key Value), which uniquely combines the training efficiency of transformers with the inference efficiency of RNNs through a dual formulation that can be expressed equivalently as either architecture. The RWKV mechanism decomposes attention-like computation into receptance (R), weight (W), key (K), and value (V) components, where receptance controls information reception, weight determines temporal importance, and key-value pairs carry information content. Unlike standard attention, RWKV employs channel-wise attention rather than position-wise attention, attending across feature dimensions rather than across sequence positions.

The dual formulation enables parallel computation during training, where the entire sequence is processed simultaneously in a transformer-like forward pass. During inference, the same computation can be reformulated as a recurrence, processing one token at a time while maintaining a constant-size hidden state. This duality provides parallelizable training with full GPU utilization while enabling constant-space, RNN-like inference without KV cache overhead. The approach achieves $O(n)$ time complexity during training with fully parallel computation, and $O(n)$ total complexity during generation with $O(1)$ space per step.

RWKV incorporates a learned time decay mechanism that determines how quickly information from previous positions fades. The exponential decay provides inductive bias toward more recent context while maintaining the ability to preserve information across long distances when relevant. This time decay mechanism serves a role analogous to positional encoding in transformers, providing temporal structure while remaining parameter-efficient.

The architecture was scaled to 14B parameters, representing the largest dense RNN trained to date and demonstrating that recurrent architectures can reach scales previously thought to require attention mechanisms. RWKV-14B achieved performance comparable to similarly-sized transformers on language modeling benchmarks while offering substantial inference efficiency advantages. The constant-memory inference proved particularly valuable for long-form generation tasks, enabling generation of arbitrarily long sequences without memory growth.

Beyond language, RWKV demonstrated versatility across multiple domains including computer vision, audio processing, and music generation. The architecture's success validated that alternatives to both attention and state space models remain viable, and that RNN-style recurrence can compete with modern architectures when designed with appropriate mechanisms for information flow and temporal modeling. The community-driven development model, with open-source implementations and collaborative improvement, established RWKV as a grassroots alternative to corporate-developed architectures.

\subsubsection{RetNet}

\citet{sun2023retentive} introduced RetNet (Retentive Network), which supports three distinct computation paradigms through a unified retention mechanism: parallel representation for training, recurrent representation for inference, and chunkwise recurrent representation for efficient long-sequence processing. The retention mechanism replaces standard attention with an exponentially-decaying similarity function that enables all three formulations through mathematical equivalence.

The parallel representation expresses retention as a matrix operation similar to attention, enabling full GPU utilization during training. The computation processes all positions simultaneously, calculating weighted combinations of values based on exponentially-decayed position-dependent weights. This parallel form supports efficient backpropagation and gradient computation, enabling RetNet to match transformer training efficiency.

The recurrent representation reformulates the identical computation as an RNN-like recurrence, where the hidden state $h_t$ evolves according to $h_t = \gamma h_{t-1} + k_t v_t^T$, with $\gamma$ representing the decay factor. This formulation enables $O(1)$ complexity per token during inference, as each step requires only updating the fixed-size hidden state with the current token's contribution. The recurrent formulation eliminates the KV cache entirely, as all history is compressed into the state vector.

The chunkwise recurrent representation divides long sequences into chunks, processes each chunk in parallel, and combines chunk outputs recurrently. Within each chunk, the parallel formulation enables efficient GPU computation. Across chunks, the recurrent formulation maintains a summary state with constant memory. This hybrid approach provides linear $O(n)$ complexity with favorable constants, enabling efficient processing of sequences that exceed GPU memory capacity for full parallel processing.

RetNet demonstrated impressive efficiency gains in practice. A 7B parameter RetNet model achieved 8.4x faster decoding than an equivalent transformer with KV cache when processing 8K token sequences. Memory consumption dropped by 70\% compared to the transformer baseline, enabled by the absence of growing KV cache. The speedup increased for longer sequences, with even more dramatic improvements at 16K or 32K context lengths. Even compared to FlashAttention-optimized transformers, RetNet maintained 7x acceleration, demonstrating that the efficiency gains stem from algorithmic advantages rather than implementation details.

The retention mechanism proved more effective than other linear attention variants, outperforming RWKV, H3, and Hyena on language modeling benchmarks. The exponential decay formulation provides better position encoding than simple linear attention, as it naturally captures the intuition that nearby tokens are typically more relevant while maintaining the ability to preserve information when needed. On downstream tasks including question answering, summarization, and classification, RetNet achieved quality comparable to standard transformers while maintaining the substantial efficiency benefits.

The three computation paradigms provide flexibility for different deployment scenarios. The parallel mode enables efficient training on GPU clusters, the recurrent mode optimizes for streaming or real-time inference, and the chunkwise mode handles batch processing of long documents. This flexibility makes RetNet suitable for diverse production environments with varying computational constraints and latency requirements.

\subsection{Hardware-Aware Optimizations}

\subsubsection{FlashAttention}

\citet{dao2022flashattention} introduced FlashAttention, representing a fundamentally different approach to efficient attention that operates orthogonally to sparsity and approximation methods. Rather than changing the attention mechanism itself, FlashAttention optimizes the implementation to account for GPU memory hierarchy, achieving dramatic speedups for exact attention computation with no approximation or quality loss. The core insight is that standard attention implementations are bottlenecked by memory bandwidth rather than arithmetic throughput, as they repeatedly read and write the attention matrix between slow high-bandwidth memory (HBM) and fast on-chip SRAM.

FlashAttention employs a tiling strategy that divides the attention computation into blocks small enough to fit entirely in fast SRAM. The algorithm loads blocks of queries, keys, and values from HBM to SRAM, computes attention within the block, and incrementally updates the output in HBM. By processing attention in tiles and never materializing the full $n \times n$ attention matrix, FlashAttention reduces memory traffic by orders of magnitude. The tiled computation requires careful handling of the softmax normalization to maintain numerical stability and exact equivalence with standard attention.

The backward pass employs a recomputation strategy, recalculating attention values during backpropagation rather than storing them. This approach trades computation for memory, an advantageous trade-off given that modern GPUs are often memory-bound rather than compute-bound. The recomputation maintains numerical stability through careful ordering of operations and use of log-space computations where appropriate. The resulting algorithm is provably IO-optimal for a range of SRAM sizes, meaning no algorithm can achieve better memory complexity while computing exact attention.

FlashAttention achieved substantial practical speedups across diverse scenarios. For BERT-large with 512-token sequences, the approach yielded 15\% end-to-end training speedup. For GPT-2 with 1K sequences, speedup reached 3x, with even greater benefits for longer sequences. On long-range tasks with 4K tokens, FlashAttention achieved 2.4x speedup while maintaining identical model quality. The memory savings proved equally impressive: 5x reduction at sequence length 512, scaling to 10x at 2K tokens and 20x at 4K tokens. These memory savings enabled longer sequences on fixed hardware, effectively expanding the accessible context length.

The IO-aware approach proved broadly applicable, benefiting any attention computation regardless of the specific attention variant. FlashAttention can be applied to sparse attention patterns, multi-head attention, cross-attention in encoder-decoder models, and various forms of local or global attention. This generality led to rapid adoption across the field, with FlashAttention becoming the de facto implementation for attention in many production systems. The approach demonstrated that algorithmic optimizations informed by hardware characteristics can yield practical benefits comparable to or exceeding mathematical approximations.

\subsubsection{FlashAttention-2}

\citet{dao2023flashattention2} introduced FlashAttention-2, addressing remaining inefficiencies in the original FlashAttention algorithm. While FlashAttention achieved 25-40\% of theoretical peak FLOPs utilization on A100 GPUs, analysis revealed that thread block distribution, shared memory communication, and non-matrix-multiply operations limited efficiency. FlashAttention-2 implemented three key optimizations: reducing non-GEMM FLOPs through algorithmic improvements, parallelizing computation across thread blocks even within a single attention head, and optimizing work distribution between warps to minimize shared memory communication.

The improvements yielded 2x speedup over FlashAttention and up to 9x speedup over standard PyTorch attention. FlashAttention-2 achieved 50-73\% of theoretical maximum FLOPs per second on A100 GPUs, approaching the efficiency of highly optimized GEMM (matrix multiplication) operations. On A100 GPUs, the implementation reached 225 TFLOPs/s with 72\% model FLOPs utilization, representing state-of-the-art efficiency for attention computation. The improved GPU occupancy through better parallelism meant that FlashAttention-2 scaled more effectively to different sequence lengths and batch sizes.

The practical impact extended beyond raw speed improvements. By nearly doubling attention efficiency, FlashAttention-2 enabled longer context windows within fixed time budgets, larger batch sizes within fixed memory budgets, and reduced training costs for large language models. The implementation maintained the exact computation guarantee, producing bitwise identical results to standard attention while operating dramatically faster. This combination of speed and exactness made FlashAttention-2 the preferred implementation for production systems where both efficiency and reproducibility matter.

Subsequent development produced FlashAttention-3, targeting newer H100 GPUs with additional optimizations for asynchronous memory operations and low-precision computation including FP8 support. The continued evolution of hardware-aware attention algorithms demonstrates that co-design of algorithms and hardware represents a sustainable path to efficiency gains, complementing mathematical improvements to attention mechanisms themselves.

\subsection{Additional Efficient Mechanisms}

\paragraph{Sparse Transformers}

Prior to Longformer and BigBird, \citet{child2019generating} introduced Sparse Transformers with factorized sparse attention patterns achieving $O(n\sqrt{n})$ complexity. The approach employed two-dimensional attention factorization, where separate attention heads attend to either rows or columns of a two-dimensional arrangement of tokens. This factorization proved particularly natural for image generation, where pixels are arranged in a spatial grid. Sparse Transformers set state-of-the-art results on density modeling for Enwik8 and demonstrated strong performance on CIFAR-10 and ImageNet 64x64 generation, establishing early evidence that sparsity could maintain quality while improving efficiency.

\paragraph{Synthesizer}

\citet{tay2020synthesizer} questioned the necessity of token-token interactions in attention through the Synthesizer architecture, which learns synthetic attention weights without query-key dot products. The Dense Synthesizer variant learns position-dependent attention patterns through feed-forward networks, while the Random Synthesizer employs fixed random attention patterns. Surprisingly, these simplified mechanisms achieved performance comparable to standard transformers on multiple benchmarks, suggesting that positional attention patterns capture much of attention's value. The Synthesizer achieved 60\% speedup over dynamic convolutions and demonstrated that the dot-product attention formulation, while effective, is not the only viable approach to sequence mixing.

\paragraph{Grouped Query Attention}

\citet{ainslie2023gqa} introduced Grouped Query Attention (GQA), an intermediate approach between Multi-Head Attention (MHA) with separate key-value pairs per head and Multi-Query Attention (MQA) with a single shared key-value pair. GQA groups queries to share key-value heads, reducing KV cache size by 10-100x while maintaining quality closer to full MHA than MQA achieves. The approach enables 12x faster decoder inference through reduced memory traffic, providing a practical middle ground between quality and efficiency. GQA has seen widespread adoption in recent large language models including Llama 2 and Mistral, validating its effectiveness in production systems.

\paragraph{Attention with Linear Biases (ALiBi)}

\citet{press2022train} proposed Attention with Linear Biases (ALiBi), replacing positional embeddings with linear biases added to attention scores proportional to the distance between tokens. This simple modification enables remarkable extrapolation properties: models trained on sequences of length $L$ can effectively process sequences of length $2L$ to $8L$ at test time. ALiBi maintains performance at 10,000+ token contexts when trained on far shorter sequences, providing a parameter-efficient alternative to learned positional embeddings. The approach has been adopted in models including MPT and BLOOM, demonstrating practical value for scenarios where test-time sequences may exceed training lengths.

\paragraph{Nyströmformer}

\citet{xiong2021nystromformer} applied the Nyström method to approximate self-attention through sampling landmark tokens. By selecting $m$ representative tokens and computing attention interactions only for this subset, the approach reconstructs approximate full attention from the sampled structure. Nyströmformer achieves $O(n)$ complexity while maintaining dense attention patterns, providing an alternative approximation strategy to kernel methods. The approach proved effective on Long Range Arena benchmarks while offering interpretability through the explicitly selected landmark tokens.

\paragraph{Hyena}

\citet{poli2023hyena} introduced Hyena, replacing attention with implicitly parameterized long convolutions and data-controlled gating. The architecture achieves subquadratic $O(n \log n)$ complexity through efficient FFT-based convolution while maintaining expressiveness through element-wise gating that depends on input content. Hyena demonstrated 20\% compute reduction at sequence length 2K compared to transformers, with speedups increasing to 2x at 8K tokens and 100x at 64K tokens. The convolutional approach provides an alternative paradigm to both attention and state space models, validating that multiple architectural families can achieve competitive efficiency and quality.

\subsection{Comparative Analysis}

The landscape of efficient attention mechanisms reveals several clear trade-offs between computational complexity, memory requirements, implementation complexity, and model quality. Fixed sparse patterns including Longformer and BigBird achieve linear complexity through structured sparsity while maintaining exact computation on the sparse support. These methods offer strong interpretability and preserve the attention paradigm, but the fixed patterns may miss task-relevant dependencies that fall outside the sparse structure. The approaches integrate easily into existing transformer codebases and benefit from theoretical guarantees of expressiveness, particularly in the case of BigBird's universal approximation results.

Linear approximations including Performer and Linformer reduce complexity through mathematical reformulation rather than sparsity. Performer's kernel method provides unbiased approximation with theoretical convergence guarantees, while Linformer's low-rank factorization exploits empirically observed structure in attention matrices. These methods achieve true $O(n)$ complexity and handle arbitrary-length sequences, but introduce approximation error that can manifest as small quality degradation on downstream tasks. The approximation error typically remains small, on the order of 1-2\% on standard metrics, but may be more significant for tasks requiring precise attention to rare events or distant dependencies.

State space models including S4 and Mamba replace attention entirely with alternative sequence-mixing operations based on continuous-time dynamics. These architectures achieve linear or log-linear training complexity and constant-memory inference, providing the strongest efficiency guarantees particularly for very long sequences. The constant-memory inference eliminates KV cache overhead, enabling generation of arbitrarily long sequences without memory growth. However, state space models employ fundamentally different mechanisms from attention, requiring different training procedures and achieving different inductive biases. The ecosystem of pretrained models and fine-tuning recipes remains less mature than for transformers, though Mamba's success has rapidly expanded available resources.

Linear-time alternatives including RWKV and RetNet provide hybrid approaches that combine transformer-like training with RNN-like inference. These architectures achieve $O(n)$ complexity with constant-space inference through dual formulations that express the same computation as either parallel or recurrent. The approaches offer practical benefits for deployment scenarios where memory efficiency matters, while maintaining training efficiency through parallelization. The retention and receptance mechanisms provide different inductive biases than attention, which can be advantageous for certain tasks but may require adaptation of training procedures developed for transformers.

Hardware-aware optimizations including FlashAttention operate orthogonally to the above categories, making any attention implementation dramatically more efficient through IO-awareness. FlashAttention maintains exact attention computation with no approximation while achieving speedups of 2-9x and memory reductions of 5-20x depending on sequence length. The approach demonstrates that algorithm-hardware co-design can yield benefits comparable to mathematical approximation without sacrificing quality. The exact computation and drop-in compatibility have led to widespread adoption, with FlashAttention becoming standard in many transformer implementations.

Hybrid architectures including Jamba demonstrate that multiple paradigms can be combined within a single model, selecting architectural components based on their strengths for specific computational roles. By interleaving attention, state space, and mixture-of-experts layers, hybrid designs optimize the efficiency-expressiveness trade-off at a granular level. This direction suggests that future models may increasingly employ heterogeneous architectures rather than committing to a single mechanism, enabling fine-grained optimization based on layer-specific requirements.

\subsection{Evaluation Metrics and Benchmarks}

The Long Range Arena benchmark established by \citet{tay2021long} provides standardized evaluation across multiple tasks requiring long-range reasoning, including ListOps for hierarchical structure, Text Classification for document-level understanding, Retrieval for similarity matching, Image Classification for sequential pixel processing, and Path-X for extreme-length path tracing. Performance on Long Range Arena has become standard for validating efficient attention mechanisms, enabling direct comparison across sparse patterns, approximations, and alternative architectures.

Beyond task-specific accuracy, efficient attention mechanisms must be evaluated on computational metrics including wall-clock training time, wall-clock inference latency, memory consumption during training, memory consumption during inference, throughput measured in tokens per second, and maximum sequence length achievable on fixed hardware. These metrics often trade off against each other—for instance, sparse patterns may increase computational FLOPs relative to their memory savings, while approximations may provide better throughput than sparse methods on modern hardware optimized for dense computation.

The practical value of efficient attention depends critically on deployment context. For offline batch processing of documents, total throughput matters most, favoring methods like FlashAttention that maximize GPU utilization. For interactive applications, per-token latency proves critical, favoring constant-memory inference methods like Mamba and RetNet. For edge deployment, total memory footprint dominates, again favoring constant-memory approaches. For research and experimentation, implementation simplicity and compatibility with existing codebases becomes paramount, favoring drop-in replacements like Longformer and FlashAttention.

\subsection{Future Directions}

Several open research questions remain in the design of efficient attention mechanisms. The optimal mixing of attention and alternative architectures in hybrid models remains unclear, with design choices currently made through empirical exploration rather than principled understanding. Better position encoding schemes for linear attention methods could improve their effectiveness on tasks requiring precise positional reasoning. Understanding how efficiency methods affect scaling laws relative to standard transformers would inform decisions about when to employ efficient mechanisms versus simply scaling compute.

Training stability for state space models and linear methods at very large scale continues to present challenges, with numerical issues and optimization difficulties that do not affect standard transformers. Learned or adaptive sparsity patterns that adjust based on task and input could outperform fixed patterns while maintaining efficiency. Hardware co-design tailored to efficient attention mechanisms, rather than retrofitting algorithms to hardware designed for dense matrix multiplication, could unlock further performance gains.

The theoretical understanding of when approximations preserve task-relevant information remains incomplete. While universal approximation results provide existence guarantees, practical understanding of the relationship between approximation error and downstream task performance would guide design decisions. Multi-modal extensions that apply efficient attention to vision, audio, and video modalities require domain-specific adaptations of the core mechanisms.

The field continues rapid evolution, with new architectures and optimizations appearing regularly. The diversity of successful approaches—sparse patterns, kernel methods, state spaces, recurrences, and hardware optimizations—suggests that no single solution dominates across all contexts. Rather, the choice of efficient attention mechanism depends on the specific requirements of sequence length, available compute, deployment constraints, and task characteristics. As language models continue to scale and expand to longer contexts, efficient attention mechanisms will remain critical to enabling progress while managing computational costs.
